% !TEX TS-program = pdflatex
% !TEX encoding = UTF-8 Unicode


%%%%%%%%%%%%%%%%%%%%%%%%%%%%%%%%%%%%%%%%%%%%%%%%%%%%%%%%%%%%%%%%%%%%%%%%%%%%
% 																							                           %
%					                         PREAMBLE      							             %
%	                                                                         %
%%%%%%%%%%%%%%%%%%%%%%%%%%%%%%%%%%%%%%%%%%%%%%%%%%%%%%%%%%%%%%%%%%%%%%%%%%%%
\documentclass[compress]{beamer}

\mode<presentation>
{
  \usetheme{Frankfurt}
  \usecolortheme{crane}
  \setbeamercovered{transparent}
}

% Package Setup
%%%%%%%%%%%%%%%%%%%%%%%%%%%%%%%%%%%%%%%%%%%%%%%%%%%%%%%%%%%%%%%%%%%%%%%%%%
%                                                                        %
%                                 PREAMBLE                               %
%                                                                        %
%%%%%%%%%%%%%%%%%%%%%%%%%%%%%%%%%%%%%%%%%%%%%%%%%%%%%%%%%%%%%%%%%%%%%%%%%%

%% PACKAGES
\usepackage[]{lineno}
\usepackage{fancyvrb}
%\linenumbers
\usepackage{amsmath}
\usepackage{microtype}
\usepackage{algorithmic}

%% GRAPHICS RELATED
\usepackage{graphicx}
\usepackage[outdir=./tmp/]{epstopdf}
\graphicspath{{../images/}{./}{./tmp/}}
\DeclareGraphicsExtensions{.eps, .pdf, .jpeg, .png}

%% CAPTION SETUP
\usepackage{float}
\usepackage[font=footnotesize]{caption}
\usepackage[font=small]{subcaption}
\captionsetup{belowskip=12pt,aboveskip=4pt}

%% BEAMER
\usepackage{multicol}
\usepackage{multirow}
\usepackage{array}				% Table Stuff
\usepackage{arydshln}
\usepackage{rotating}

%% BIBLIOGRAPHY
\bibliographystyle{ieeetr}

%% UNITS
\usepackage{siunitx}

%% EQUATIONS
%\numberwithin{equation}{section}

%%%%%%%%%%%%%%%%%%%%%%%%%%%%%%%%%%%%%%%%%%%%%%%%%%%%%%%%%%%%%%%%%%%%%%%%%%%
%                                                                         %
%                             Listing Setup                               %
%                                                                         %
%%%%%%%%%%%%%%%%%%%%%%%%%%%%%%%%%%%%%%%%%%%%%%%%%%%%%%%%%%%%%%%%%%%%%%%%%%%
\usepackage{listings}
\lstset{ %
    language=C++,
    basicstyle=\footnotesize\ttfamily,
    numbers=left,
    numberstyle=\tiny\color{gray},
    stepnumber=2,
    numbersep=5pt,
    backgroundcolor=\color{white},
    showspaces=false,
    showstringspaces=false,
    showtabs=false,
    frame=single,
    rulecolor=\color{black},
    tabsize=2,
    breaklines=true,
    breakatwhitespace=false,
    title=\lstname,
    keywordstyle=\color{blue},
    commentstyle=\color{OliveGreen},
    stringstyle=\color{orange}
}
\DeclareCaptionFont{white}{\color{white}}
\DeclareCaptionFormat{listing}{\colorbox[cmyk]{0.43, 0.35, 0.35, 0.01}{\parbox{\dimexpr\textwidth-2\fboxsep\relax}{#1#2#3}}}
\captionsetup[lstlisting]{format=listing,labelfont=white,textfont=white,singlelinecheck=false,margin=0pt,font={bf,footnotesize}}
%\lstnewenvironment{code}[1][]%
%{ \noindent\minipage{\linewidth}
%	\lstset{#1}
%}
%{\endminipage}
%% USER COMMANDS
\usepackage{isotope}
\newcommand{\iso}{\isotope}
\newcommand{\figurewidth}{\textwidth}
\newcommand{\micron}{$\mu$m}



% Preamble / Frst Size
%\setbeamersize{text margin left=5mm, text margin right 5mm}
\title[RPM8 Optimiziation] {Optimization of the RPM8 Detector with Genetic Algorthims}
\author[] {
    Matthew Urffer\inst{1}
}
\institute[University of Tennessee] { 
  \inst{1}%
  Department of Nuclear Engineering,
  University of Tennessee, Knoxville, TN

}

\date[] {March 3, 2013}
\pgfdeclareimage[height=0.5cm]{university-logo}{images/utwordmarkhorz.eps}
\logo{\pgfuseimage{university-logo}}

% Delete this, if you do not want the table of contents to pop up at
% the beginning of each subsection:
\AtBeginSection[]
{
  \begin{frame}<beamer>{}
    \frametitle{\insertsectionhead}
    \begin{multicols}{2}
    \tableofcontents[currentsection] 
    \end{multicols}
  \end{frame}
}

\begin{document}

\begin{frame}[plain]
  \titlepage
  \tiny
    \begin{center}
\centering{Financial support from the Domestic Nuclear Detection Office (DNDO) through Award No. 003387891 is gratefully acknowledged. 
  Any opinions, findings, and conclusions or recomendations expressed in this material are those of the presenter and do not necessarily reflect the views of DNDO.}
  \end{center}
\end{frame}

\begin{frame}{Table of Contents}
  \begin{multicols}{2}
    %\tableofcontents[currentsection,pausesections]
    \tableofcontents[currentsection]
    \end{multicols}
\end{frame}


%%%%%%%%%%%%%%%%%%%%%%%%%%%%%%%%%%%%%%%%%%%%%%%%%%%%%%%%%%%%%%%%%%%%%%%%%%%%
%                                                                          %
%                            START OF CONTENT                              %
%                                                                          %
%%%%%%%%%%%%%%%%%%%%%%%%%%%%%%%%%%%%%%%%%%%%%%%%%%%%%%%%%%%%%%%%%%%%%%%%%%%%
\include{./PyevovleIntroPresentation.tex}
\section*{Conclusions}
%%%%%%%%%%%%%%%%%%%%%%%%%%%%%%%%%%%%%%%%%%%%%%%%%%%%%%%%%%%%%%%%%%%%%%%%%%%%
\begin{frame}{Summary}

  \begin{itemize}
  \small
  \item
    A framework has been developed for the characterization of possible replacement technologies for radiation portal monitors\cite{NotHere}
    \begin{itemize}
	        \tiny
			\item Verification of instrument gains
			\item Neutron Measurements
			\item Gamma Measurements
			\item Various button sources
		\end{itemize}
    \small
		\item Ability to calculate MLLD necessary for $\epsilon_{int,\gamma} \le 10^{-6}$
    A framework has been developed for pulse shape discrimination 
  \item
    Thin polymeric films have been demonstrated to have the necessary interaction rates for radiation portal monitors
  \end{itemize}
\begin{figure}
	\centering
		\includegraphics[width=0.2\textwidth]{Questions.eps}
\end{figure}
\end{frame}

%%%%%%%%%%%%%%%%%%%%%%%%%%%%%%%%%%%%%%%%%%%%%%%%%%%%%%%%%%%%%%%%%%%%%%%%%%%
% BILBIOLGRAPHY
\begin{frame}[plain,allowframebreaks]
\frametitle{Works Cited}
	\tiny
  \bibliography{../Zotero}
\end{frame}

%%%%%%%%%%%%%%%%%%%%%%%%%%%%%%%%%%%%%%%%%%%%%%%%%%%%%%%%%%%%%%%%%%%%%%%%%%%
% APPENDIX

\end{document}


